% Options for packages loaded elsewhere
\PassOptionsToPackage{unicode}{hyperref}
\PassOptionsToPackage{hyphens}{url}
\documentclass[
]{article}
\usepackage{xcolor}
\usepackage[margin=1in]{geometry}
\usepackage{amsmath,amssymb}
\setcounter{secnumdepth}{-\maxdimen} % remove section numbering
\usepackage{iftex}
\ifPDFTeX
  \usepackage[T1]{fontenc}
  \usepackage[utf8]{inputenc}
  \usepackage{textcomp} % provide euro and other symbols
\else % if luatex or xetex
  \usepackage{unicode-math} % this also loads fontspec
  \defaultfontfeatures{Scale=MatchLowercase}
  \defaultfontfeatures[\rmfamily]{Ligatures=TeX,Scale=1}
\fi
\usepackage{lmodern}
\ifPDFTeX\else
  % xetex/luatex font selection
\fi
% Use upquote if available, for straight quotes in verbatim environments
\IfFileExists{upquote.sty}{\usepackage{upquote}}{}
\IfFileExists{microtype.sty}{% use microtype if available
  \usepackage[]{microtype}
  \UseMicrotypeSet[protrusion]{basicmath} % disable protrusion for tt fonts
}{}
\makeatletter
\@ifundefined{KOMAClassName}{% if non-KOMA class
  \IfFileExists{parskip.sty}{%
    \usepackage{parskip}
  }{% else
    \setlength{\parindent}{0pt}
    \setlength{\parskip}{6pt plus 2pt minus 1pt}}
}{% if KOMA class
  \KOMAoptions{parskip=half}}
\makeatother
\usepackage{graphicx}
\makeatletter
\newsavebox\pandoc@box
\newcommand*\pandocbounded[1]{% scales image to fit in text height/width
  \sbox\pandoc@box{#1}%
  \Gscale@div\@tempa{\textheight}{\dimexpr\ht\pandoc@box+\dp\pandoc@box\relax}%
  \Gscale@div\@tempb{\linewidth}{\wd\pandoc@box}%
  \ifdim\@tempb\p@<\@tempa\p@\let\@tempa\@tempb\fi% select the smaller of both
  \ifdim\@tempa\p@<\p@\scalebox{\@tempa}{\usebox\pandoc@box}%
  \else\usebox{\pandoc@box}%
  \fi%
}
% Set default figure placement to htbp
\def\fps@figure{htbp}
\makeatother
\setlength{\emergencystretch}{3em} % prevent overfull lines
\providecommand{\tightlist}{%
  \setlength{\itemsep}{0pt}\setlength{\parskip}{0pt}}
\usepackage{bookmark}
\IfFileExists{xurl.sty}{\usepackage{xurl}}{} % add URL line breaks if available
\urlstyle{same}
\hypersetup{
  pdftitle={index},
  pdfauthor={Alex Medina},
  hidelinks,
  pdfcreator={LaTeX via pandoc}}

\title{index}
\author{Alex Medina}
\date{2026-10-18}

\begin{document}
\maketitle

\section{PART II}\label{part-ii}

\subsection{(2.1.1) Scrape a Table}\label{scrape-a-table}

table\_url \textless-
``\url{https://en.wikipedia.org/wiki/Community_areas_of_Chicago}''

table\_page \textless- read\_html(table\_url)

\subsubsection{COMMUNITY AREA AND TOTAL
POPULATION}\label{community-area-and-total-population}

wiki\_tables \textless- table\_page \%\textgreater\%
html\_elements(``table.wikitable'') \%\textgreater\% html\_table()

population\_table \textless- wiki\_tables{[}{[}2{]}{]}

\subsection{(2.1.2)}\label{section}

population\_table \textless- population\_table \%\textgreater\%
clean\_names() \%\textgreater\% filter(name != ``Chicago'')
\%\textgreater\% mutate(total\_population =
as.numeric(str\_remove\_all(total\_population, ``,'')))

\subsection{(2.1.3)}\label{section-1}

population\_table \%\textgreater\% filter(name == ``West Town'')

population\_table \%\textgreater\% arrange(desc(total\_population))
\%\textgreater\% head(5)

nrow(population\_table) head(population\_table)

\subsection{(2.2.1) Scrape with CSS
Selectors}\label{scrape-with-css-selectors}

css\_url = ``\url{https://books.toscrape.com/}'' css\_page =
read\_html(css\_url)

\subsubsection{Title, Price,
Availabiliity}\label{title-price-availabiliity}

\subsubsection{Title: article.product\_prod -\textgreater{} h3
-\textgreater{} a -\textgreater{}
title}\label{title-article.product_prod---h3---a---title}

\subsubsection{Price: article.product\_prod -\textgreater{}
div.product\_price -\textgreater{}
p.price\_color}\label{price-article.product_prod---div.product_price---p.price_color}

\subsubsection{Availability: article.product\_prod -\textgreater{}
div.product\_price -\textgreater{}
p.instock.availability}\label{availability-article.product_prod---div.product_price---p.instock.availability}

books\_table \textless- css\_page \%\textgreater\%
html\_elements(``article.product\_pod'')

title \textless- books\_table \%\textgreater\% html\_element(``h3 a'')
\%\textgreater\% html\_attr(``title'')

price \textless- books\_table \%\textgreater\%
html\_element(``p.price\_color'') \%\textgreater\% html\_text2()

availability \textless- books\_table \%\textgreater\%
html\_element(``p.instock.availability'') \%\textgreater\% html\_text2()

books \textless- tibble( title, price, availability )

\subsection{(2.2.2)}\label{section-2}

books \textless- books \%\textgreater\% mutate(price =
as.numeric(str\_remove\_all(price, ``£'')))

write\_csv(books, ``books.csv'')

nrow(books) head(books)

\subsection{(2.2.BONUS)}\label{bonus}

\subsubsection{star-rating (\# of stars): article.product\_pod
-\textgreater{} p.star-rating.(\# of
stars)}\label{star-rating-of-stars-article.product_pod---p.star-rating.-of-stars}

\subsubsection{Grab class star-rating (\# of
stars)}\label{grab-class-star-rating-of-stars}

\subsubsection{Remove ``star-rating'' to leave \# of
stars}\label{remove-star-rating-to-leave-of-stars}

stars \textless- books\_table \%\textgreater\%
html\_element(``p.star-rating'') \%\textgreater\% html\_attr(``class'')
\%\textgreater\% str\_remove(``star-rating'')

books \textless- books \%\textgreater\% mutate(rating = ifelse(stars ==
``One'', 1, ifelse(stars == ``Two'', 2, ifelse(stars == ``Three'', 3,
ifelse(stars == ``Four'', 4, 5)))))

write\_csv(books, ``books.csv'')

nrow(books) head(books)

\section{Part III}\label{part-iii}

\subsection{(3.1)}\label{section-3}

open\_url \textless-
``\url{https://data.cityofchicago.org/resource/v6vf-nfxy.json}''
parameters \textless- list( community\_area = ``24'', \texttt{\$where} =
``created\_date \textgreater= `2026-01-01T00:00:00'\,'',
\texttt{\$order} = ``created\_date DESC'', \texttt{\$limit} = ``1000'' )

\subsection{(3.2)}\label{section-4}

data \textless- GET(open\_url, parameters)

west\_town \textless- fromJSON(content(data, ``text''), flatten = TRUE)
nrow(west\_town) names(west\_town)

\subsection{(3.3)}\label{section-5}

top\_ten \textless- west\_town \%\textgreater\% count(sr\_type, sort =
TRUE) \%\textgreater\% head(10)

top\_ten

\subsection{(3.4)}\label{section-6}

\subsubsection{The most common requests in West Town are 311 Information
Only Calls (305 calls) and Aircraft Noise Complaints (304 calls) which
suggests that majority of residents likely seek information or concerns
about about flight-path noise. What the data doesn't really tell us is
if there are only a small amount of people calling multiple times about
311/Aircraft Noise or if it is a different person every time as the
difference in calls is quite large with the third report only having 75
calls).}\label{the-most-common-requests-in-west-town-are-311-information-only-calls-305-calls-and-aircraft-noise-complaints-304-calls-which-suggests-that-majority-of-residents-likely-seek-information-or-concerns-about-about-flight-path-noise.-what-the-data-doesnt-really-tell-us-is-if-there-are-only-a-small-amount-of-people-calling-multiple-times-about-311aircraft-noise-or-if-it-is-a-different-person-every-time-as-the-difference-in-calls-is-quite-large-with-the-third-report-only-having-75-calls.}

\subsection{(3.5)}\label{section-7}

\section{FINISH IF YOU HAVE TIME}\label{finish-if-you-have-time}

top\_five \textless- west\_town \%\textgreater\% count(sr\_type, sort =
TRUE) \%\textgreater\% head(5) \%\textgreater\% pull(sr\_type)

\subsubsection{close\_days \textless- west\_town
\%\textgreater\%}\label{close_days---west_town}

\subsubsection{filter(status == ``Completed'', sr\_type \%in\%
top\_five)}\label{filterstatus-completed-sr_type-in-top_five}

\section{PART IV}\label{part-iv}

\subsection{(4.1)}\label{section-8}

\subsubsection{Gross Rent as a Percentage of Household Income,
``Total'':
B25070\_001E}\label{gross-rent-as-a-percentage-of-household-income-total-b25070_001e}

\subsubsection{• Gross Rent as a Percentage of Household Income, ``50.0
percent or more'': \_\_\_\_\_\_\_\_
B25070\_010E}\label{gross-rent-as-a-percentage-of-household-income-50.0-percent-or-more-________-b25070_010e}

\subsubsection{• Median Household Income: \_\_\_\_\_\_\_\_
B19013\_001E}\label{median-household-income-________-b19013_001e}

\subsection{(4.2)}\label{section-9}

ohio \textless- get\_acs( geography = ``county'', variables = c( Total =
``B25070\_001E'', Rent\_50 = ``B25070\_010E'', median\_income =
``B19013\_001E'' ), state = ``OH'', year = 2024, survey = ``acs5'' )

\subsection{(4.3)}\label{section-10}

ohio\_wider \textless- ohio \%\textgreater\% select(GEOID, NAME,
variable, estimate) \%\textgreater\% pivot\_wider(names\_from =
variable, values\_from = estimate)

ohio\_wider \textless- ohio\_wider \%\textgreater\%
mutate(severe\_burden\_rate = B25070\_010 / B25070\_001 * 100)
\%\textgreater\% arrange(desc(severe\_burden\_rate))

\subsection{(4.BONUS)}\label{bonus-1}

ohio\_wider \textless- ohio\_wider \%\textgreater\% separate(NAME, into
= c(``county'', ``state''), ``,'')

\subsection{(4.4)}\label{section-11}

ohio\_wider \%\textgreater\% summarize( mean\_rate =
mean(severe\_burden\_rate), median\_rate = median(severe\_burden\_rate)
)

top5 \textless- ohio\_wider \%\textgreater\%
arrange(desc(severe\_burden\_rate)) \%\textgreater\% head(5)

bottom5 \textless- ohio\_wider \%\textgreater\%
arrange(severe\_burden\_rate) \%\textgreater\% head(5)

\subsection{(4.5.1)}\label{section-12}

ggplot(ohio\_wider, aes(x = B19013\_001, y = severe\_burden\_rate)) +
geom\_point() + geom\_smooth(method = ``lm'') + labs( title = ``Median
Household Income vs Severely Rent Burdened Renters in Ohio (ACS 2024)'',
x = ``Median Household Income (\$)'', y = ``Severely Rent-Burdened
Renters (\%)'' )

\subsection{(4.5.2)}\label{section-13}

bar \textless- bind\_rows(top5, bottom5)

ggplot(bar, aes(x = reorder(county, severe\_burden\_rate), y =
severe\_burden\_rate)) + geom\_col() + coord\_flip() + labs(title =
``Ohio Counties with Highest and Lowest Severe Rent Burden (ACS 2024)'',
x = ``County'', y = ``Severely Rent-Burdened Renters (\%)'' )

\subsection{(4.6.1)}\label{section-14}

\section{There isn't really a strong correlation between income and
severe rent burden in Ohio. The line does slope downward, so
higher-income counties tend to have a lower severe rent burden. However,
the points are very widely scattered so income alone does not likely
explain much of the variation. I think rent burdening could apply more
to specific neighborhoods or areas rather than household
income.}\label{there-isnt-really-a-strong-correlation-between-income-and-severe-rent-burden-in-ohio.-the-line-does-slope-downward-so-higher-income-counties-tend-to-have-a-lower-severe-rent-burden.-however-the-points-are-very-widely-scattered-so-income-alone-does-not-likely-explain-much-of-the-variation.-i-think-rent-burdening-could-apply-more-to-specific-neighborhoods-or-areas-rather-than-household-income.}

\subsection{(4.6.2)}\label{section-15}

\section{One limitation could probably be how populated the county is.
After looking it up, Athens County and Portage County both have
universities and college students are more likely have lower income yet
pay some kind of rent whether that is through university housing or
housing nearby campus. Another limitation could be the amount of places
available to rent in different counties compared to just purchasing
houses.}\label{one-limitation-could-probably-be-how-populated-the-county-is.-after-looking-it-up-athens-county-and-portage-county-both-have-universities-and-college-students-are-more-likely-have-lower-income-yet-pay-some-kind-of-rent-whether-that-is-through-university-housing-or-housing-nearby-campus.-another-limitation-could-be-the-amount-of-places-available-to-rent-in-different-counties-compared-to-just-purchasing-houses.}

\section{PART V}\label{part-v}

\subsection{(5.1)}\label{section-16}

\section{Publishing exact addresses and coordinates of 311 requests
could most definitely identify the people who made them. They had put a
request for information, not to have them exposed. For example, with the
aircraft noise complaints, having the address exposed could lead to some
kind of retaliation if the complaints lead into some kind of trouble. I
would likely use aggregation in order to void the street address and
exact coordinates entirely. I would likely just leave it as the
neighborhood. I would also suppress the number of requests that is shown
if it shows to be too small. I would likely not disclose the number of
requests if they are less than
\textasciitilde5.}\label{publishing-exact-addresses-and-coordinates-of-311-requests-could-most-definitely-identify-the-people-who-made-them.-they-had-put-a-request-for-information-not-to-have-them-exposed.-for-example-with-the-aircraft-noise-complaints-having-the-address-exposed-could-lead-to-some-kind-of-retaliation-if-the-complaints-lead-into-some-kind-of-trouble.-i-would-likely-use-aggregation-in-order-to-void-the-street-address-and-exact-coordinates-entirely.-i-would-likely-just-leave-it-as-the-neighborhood.-i-would-also-suppress-the-number-of-requests-that-is-shown-if-it-shows-to-be-too-small.-i-would-likely-not-disclose-the-number-of-requests-if-they-are-less-than-5.}

\subsection{(5.2)}\label{section-17}

\section{Some people are unable to access technology, do not speak the
language that 311 may host with, and/or are just unaware of the 311
system. This would lead to many requests not being put in and the only
requests that are being noticed are biased toward English Speakers who
are able to access technology easier. A city manager would likely see
plenty of requests for one area in a primarily English Speaking town and
would first go there but would not check other areas that are not the
majority. One way to reduce the bias is to add more accessibility /
translation options with
311.}\label{some-people-are-unable-to-access-technology-do-not-speak-the-language-that-311-may-host-with-andor-are-just-unaware-of-the-311-system.-this-would-lead-to-many-requests-not-being-put-in-and-the-only-requests-that-are-being-noticed-are-biased-toward-english-speakers-who-are-able-to-access-technology-easier.-a-city-manager-would-likely-see-plenty-of-requests-for-one-area-in-a-primarily-english-speaking-town-and-would-first-go-there-but-would-not-check-other-areas-that-are-not-the-majority.-one-way-to-reduce-the-bias-is-to-add-more-accessibility-translation-options-with-311.}

\subsection{(5.3)}\label{section-18}

\section{I believe the scraping that I did in Part II was appropriate.
The Wikipedia table is public and contains no personal information.
Following its reference link, it led me to the ``Chicago Metropolitan
Agency for Planning''. The website is a government website with the data
public as well. I did nothing that would be out of the ordinary for this
information to be used nor would it violate any citizen's privacy. For
the CSS portion, ``Books to Scrape'' is a site built for scraping
practice so that also has no ethical concerns either. A situation that
would require further review would be with a dataset that holds more
private information like addresses, names, or other personal
identifiers. This would likely be anything related to court, eviction,
and other
records.}\label{i-believe-the-scraping-that-i-did-in-part-ii-was-appropriate.-the-wikipedia-table-is-public-and-contains-no-personal-information.-following-its-reference-link-it-led-me-to-the-chicago-metropolitan-agency-for-planning.-the-website-is-a-government-website-with-the-data-public-as-well.-i-did-nothing-that-would-be-out-of-the-ordinary-for-this-information-to-be-used-nor-would-it-violate-any-citizens-privacy.-for-the-css-portion-books-to-scrape-is-a-site-built-for-scraping-practice-so-that-also-has-no-ethical-concerns-either.-a-situation-that-would-require-further-review-would-be-with-a-dataset-that-holds-more-private-information-like-addresses-names-or-other-personal-identifiers.-this-would-likely-be-anything-related-to-court-eviction-and-other-records.}

\subsubsection{AI USAGE}\label{ai-usage}

\section{Used it a few times primarily for syntax/grammar errors. I
could not figure out what was wrong and realized it was someting as
simple as ``estimmate'' and I made sure my wording made sense for
explanations considering at the moment of doing the assignment I had put
a concerning amount of typos that I would not have caught. I also
checked my scatterplot code to make sure it was right because the plot
itself did not look quite right and I was
concerned.}\label{used-it-a-few-times-primarily-for-syntaxgrammar-errors.-i-could-not-figure-out-what-was-wrong-and-realized-it-was-someting-as-simple-as-estimmate-and-i-made-sure-my-wording-made-sense-for-explanations-considering-at-the-moment-of-doing-the-assignment-i-had-put-a-concerning-amount-of-typos-that-i-would-not-have-caught.-i-also-checked-my-scatterplot-code-to-make-sure-it-was-right-because-the-plot-itself-did-not-look-quite-right-and-i-was-concerned.}

\end{document}
